% DEMONSTRATION DRAFT — §2.5 Spectral Models, written in the Li et al. (2021)
% voice (present-tense "we fit", method-then-justification, grouped "e.g." cites),
% with model descriptions sourced from Ravasio et al. (2SBPL/synchrotron) and
% Guiriec et al. (multi-component thermal+non-thermal). Data-independent; ready to
% drop into the Phase-6 manuscript. %chk equation conventions vs astromodels.

\subsection{Spectral Models} \label{sec:models}
We fit each time-resolved spectrum with a family of six photon models, chosen to
span the two physical interpretations of prompt-emission curvature while
remaining empirical and directly comparable with the Fermi/GBM spectral catalogs
\citep[e.g.,][]{Kaneko2006,Gruber2014,Yu2016}. We use empirical functions
deliberately: they are flexible and model-agnostic, and although their parameters
do not map uniquely onto a radiation mechanism, the low-energy index $\alpha$ in
particular still discriminates between emission scenarios when the data are of
high significance \citep[e.g.,][]{Acuner2020}.

The first three are single-component non-thermal models. The Band function
\citep{Band1993} is a smoothly-joined broken power law with low- and high-energy
photon indices $\alpha$ and $\beta$ and a $\nu F_\nu$ peak energy \Ep; below the
break $E_b=(\alpha-\beta)\Ep/(2+\alpha)$,
\begin{equation}
N(E)=A\,(E/E_0)^{\alpha}\,\exp\!\left[-(2+\alpha)E/\Ep\right],
\end{equation}
joining $N(E)\propto E^{\beta}$ above $E_b$ ($E_0=100$~keV). The cutoff power law
(CPL) replaces the high-energy power law with an exponential cutoff,
$N(E)=A\,(E/E_0)^{\Gamma}\exp(-E/E_c)$, peaking at $\Ep=(2+\Gamma)E_c$; it is the
limit of the Band function as $\beta\!\to\!-\infty$. The smoothly broken power law
(SBPL) describes a single break with adjustable curvature.

To represent the synchrotron interpretation we add the double smoothly broken
power law \citep[2SBPL;][]{Ravasio2018,Ravasio2019}. The 2SBPL has three
power-law segments with indices $\alpha_1,\alpha_2,\beta$ joined at two break
energies: a lower break $x_b$ and an upper break $x_p$ (the latter coinciding with
the $\nu F_\nu$ peak). Under the optically-thin synchrotron interpretation the two
breaks are physical: $x_b$ is the cooling frequency \nuc\ and $x_p=\num$ is the
injection frequency, with the segment between them carrying the characteristic
slopes $\alpha_1\simeq-2/3$ (slow cooling) steepening toward $-3/2$ (fast cooling)
\citep[e.g.,][]{Sari1998,Oganesyan2017,Ravasio2019}. We hold the break
smoothnesses fixed at their default values, so the free parameters are the three
indices, the two breaks, and the normalization.

To represent the photospheric (thermal) interpretation we add a blackbody to the
two simplest non-thermal continua, following the multi-component decomposition of
\citet{Guiriec2011,Guiriec2015}, in which a sub-dominant quasi-thermal component
($\sim$Planck) accompanies the non-thermal emission and the two are fit
\emph{simultaneously}. The blackbody photon spectrum is
\begin{equation}
N(E)\propto E^{2}/\!\left[\exp(E/\kT)-1\right],
\end{equation}
with temperature \kT, and we form the composites Band$+$BB and CPL$+$BB. As in
\citet{Guiriec2015} we quantify the thermal contribution by the flux ratio
$F_{\rm BB}/F_{\rm tot}$. A BB$+$SBPL composite reproduces the curvature of a
2SBPL near the peak, so the $+$BB models act as a practical thermal proxy for the
two-break model, and comparing them is the central test of this work.

For all peaked models we report \Ep\ as the true $\nu F_\nu$ peak. The two $+$BB
models and the 2SBPL are the competing descriptions of any curvature beyond a
single spectral break --- a thermal component versus a second synchrotron break.
All fits are performed in the Multi-Mission Maximum Likelihood framework
\citep[3ML;][]{Vianello2015} with the \texttt{pgstat} likelihood (a Poisson
source with Gaussian background).
